\section{Kikuchi Bands And The Gnomonic Projection}

This note fixes the geometry that PyTex uses for Kikuchi bands, in both the EBSD
and the TEM setting. It is the theory behind
\texttt{pytex.diffraction.kikuchi}: \texttt{GnomonicProjection},
\texttt{KikuchiBand}, \texttt{KikuchiZoneAxis}, \texttt{KikuchiPattern}, and
\texttt{simulate\_kikuchi\_pattern}.

The companion user-facing pages are the EBSD and diffraction foundation concept
pages, and the executable worked examples
\texttt{diffraction-ni-111-kikuchi-band-width} and
\texttt{diffraction-gnomonic-zone-axis-radius}.

\subsection{Origin Of A Band}

Inelastic scattering inside the specimen produces electrons travelling in
essentially every direction from a small source volume. For a lattice plane of
interplanar spacing $d$, those directions making the Bragg angle $\theta_B$ with
the plane satisfy the diffraction condition, where

\[
  \sin\theta_B = \frac{\lambda}{2 d}.
\]

The set of such directions is not a single cone but two: one on each side of the
plane. Each is a cone of semi-angle $90^\circ - \theta_B$ about the plane normal
$\hat{\mathbf{n}}$, that is, the locus

\[
  \hat{\mathbf{r}} \cdot \hat{\mathbf{n}} = \pm \sin\theta_B .
\]

These are the \emph{Kossel cones}. Their intersections with the detector are the
two edges of the Kikuchi band, and the trace of the lattice plane itself,
$\hat{\mathbf{r}} \cdot \hat{\mathbf{n}} = 0$, runs midway between them.

The band therefore has \emph{angular} width exactly $2\theta_B$, independent of
detector geometry. Combined with Bragg's law this makes band width a direct
measurement of the lattice: a wide band means a large $d$-spacing.

\subsection{The Gnomonic Projection}

Let the detector plane lie at distance $L$ along the detector normal, with
in-plane orthonormal axes $\hat{\mathbf{u}}, \hat{\mathbf{v}}$ and normal
$\hat{\mathbf{n}}_d$, all expressed in the laboratory frame. The gnomonic
projection of a unit direction $\hat{\mathbf{r}}$ is the point where the ray from
the source along $\hat{\mathbf{r}}$ meets that plane, in units of $L$:

\[
  x = \frac{\hat{\mathbf{r}} \cdot \hat{\mathbf{u}}}
           {\hat{\mathbf{r}} \cdot \hat{\mathbf{n}}_d},
  \qquad
  y = \frac{\hat{\mathbf{r}} \cdot \hat{\mathbf{v}}}
           {\hat{\mathbf{r}} \cdot \hat{\mathbf{n}}_d},
  \qquad
  \hat{\mathbf{r}} \cdot \hat{\mathbf{n}}_d > 0 .
\]

Only the forward hemisphere projects; PyTex reports directions with a
non-positive denominator as invalid rather than assigning them a coordinate.

Two properties follow immediately and are the reason this projection is used.

\paragraph{Radius is the tangent of the polar angle.}
A direction at angle $\psi$ from the detector normal projects to radius
$\tan\psi$. In particular a direction $45^\circ$ from the normal lands at radius
exactly $1$. This fixes the scale of the projection without reference to any
crystallography, and is the check used in the worked example
\texttt{diffraction-gnomonic-zone-axis-radius}.

\paragraph{Great circles map to straight lines.}
Write $\hat{\mathbf{n}} = (a, b, c)$ in the detector basis. A direction with
gnomonic coordinates $(x, y)$ is proportional to $(x, y, 1)$, so the plane trace
$\hat{\mathbf{r}} \cdot \hat{\mathbf{n}} = 0$ becomes

\[
  a x + b y + c = 0,
\]

a straight line whose coefficients are simply the plane normal in the detector
basis. This holds for \emph{any} detector tilt. It is why band detection is
performed in gnomonic coordinates: in raw detector pixels the same trace is
curved whenever the detector is tilted.

\subsection{Band Edges Are Conics, Not Lines}

The Kossel cones are not great circles, so the same substitution gives

\[
  \bigl(a x + b y + c\bigr)^2 = \sin^2\theta_B \,\bigl(x^2 + y^2 + 1\bigr),
\]

a conic section --- a hyperbola for the small Bragg angles of electron
diffraction. The frequent statement that band edges are straight and parallel to
the centre line is the small-angle approximation. PyTex does not make it: edges
are obtained by sampling the exact cones

\[
  \hat{\mathbf{r}}(\phi) = \pm\sin\theta_B\, \hat{\mathbf{n}}
    + \cos\theta_B \bigl(\cos\phi\, \hat{\mathbf{e}}_1
                         + \sin\phi\, \hat{\mathbf{e}}_2\bigr),
\]

with $(\hat{\mathbf{e}}_1, \hat{\mathbf{e}}_2, \hat{\mathbf{n}})$ orthonormal,
and projecting each sampled direction. Sampling the cone rather than solving the
projected conic avoids branch bookkeeping and stays correct at any Bragg angle.

\subsection{Apparent Band Width}

The angular width $2\theta_B$ is fixed, but the width \emph{measured in gnomonic
units} is not: the projection stretches with distance from the pattern centre, so
a band far from the centre appears wider. Restricting attention to the plane
spanned by $\hat{\mathbf{n}}_d$ and the in-plane part of $\hat{\mathbf{n}}$, a
unit ray at angle $\psi$ from the detector normal has gnomonic offset $\tan\psi$
and satisfies

\[
  \hat{\mathbf{r}} \cdot \hat{\mathbf{n}}
    = q \sin\psi + c \cos\psi
    = \sin(\psi + \varphi),
  \qquad
  q = \sqrt{a^2 + b^2},
  \quad
  \varphi = \operatorname{atan2}(c, q),
\]

using $\lVert \hat{\mathbf{n}} \rVert = 1$. The Kossel condition then gives
$\psi_\pm = \pm\theta_B - \varphi$ and an apparent width

\[
  w = \bigl\lvert \tan(\theta_B - \varphi) - \tan(-\theta_B - \varphi) \bigr\rvert .
\]

This is \texttt{KikuchiBand.width\_at\_pattern\_center}. It diverges as either
edge approaches grazing incidence, where PyTex reports an infinite width rather
than a large finite number that would read as a measurement. Converting a
measured apparent width to a $d$-spacing without accounting for the band's
position is a standard source of error; exposing $w$ separately from
$2\theta_B$ makes the distinction explicit.

\subsection{Zone Axes}

A direction $[uvw]$ lies in the plane $(hkl)$ when the zone law
$hu + kv + lw = 0$ holds. Every band whose plane contains a given zone axis
therefore has a centre line passing through the projected point of that axis, so
zone axes appear as hubs where several bands intersect. PyTex enumerates
candidate axes, counts the simulated bands satisfying the zone law for each, and
reports only those shared by at least two bands, since a single band does not
define a hub.

\subsection{Frame Chain}

Plane normals are computed in the crystal frame through the reciprocal basis,
carried to the specimen frame by the crystal-to-specimen orientation matrix
$\mathbf{g}$, and then to the laboratory frame by the specimen-to-laboratory
rotation of the diffraction geometry:

\[
  \hat{\mathbf{n}}_{\mathrm{lab}}
    = \mathbf{M}_{\mathrm{spec} \rightarrow \mathrm{lab}} \,
      \mathbf{g} \, \hat{\mathbf{n}}_{\mathrm{crystal}} .
\]

This is the same chain the kinematic spot simulation uses, and it is pinned by a
closed-form test: for a cubic phase at the cube orientation with an untilted
detector, $[011]$ must project to gnomonic radius exactly $1$. A transposed or
misordered rotation anywhere in the chain moves that value.

\subsection{Current Limits}

\begin{itemize}
  \item Band positions and widths are exact for the stated geometry; intensities
        are a kinematic $\lvert F \rvert^2$ proxy.
  \item The excess/deficiency asymmetry across a band, which arises from the
        dynamical theory, is not modelled --- simulated bands are symmetric.
  \item Higher-order Laue zone rings, background, and detector point-spread are
        not modelled.
  \item Orientation determination from a measured pattern uses band
        \emph{positions}, so the intensity limitation does not affect the
        geometric use of this surface.
\end{itemize}
